\section{Surrogate-assisted local CDF estimation}
\label{sec:method}

\subsection{Target and reference-outcome benchmark}

We observe independent copies of \(O=(W,R,RY)\), where \(W\) contains \(X\),
\(S\), and any other broadly observed variables. We assume
\(R\perp Y\mid W\). In the primary designs \(W=(X,S)\); an application may
also include an observed variable used only by the verification design. The
main procedure treats the verification
probability \(\pi_M(W)\) as known, as in a designed validation sample; the
Supplement covers estimated probabilities. The reference outcome \(Y\) defines
the estimand. The surrogate \(S\) enters only through an augmentation used to
improve precision. For \(Z_y=1(Y\le y)\), let
\(q_{0,y}(X)=E(Z_y\mid X)\) and \(m_{0,y}(W)=E(Z_y\mid W)\). At an interior point
\(x_0\), we distinguish the scientific target from the target at the chosen
localization scale:
\begin{equation}
\begin{aligned}
F_0(y\mid x_0)
&=\Pr(Y\le y\mid X=x_0),\\
F_{h,0}(y\mid x_0)
&=
\frac{E\{K_h(X-x_0)Z_y\}}{\mu_h},
\qquad
\mu_h=E\{K_h(X-x_0)\},
\end{aligned}
\label{eq:regular-target}
\end{equation}
where \(K_h(X-x_0)=K\{(X-x_0)/h\}\). For a declared bandwidth, our
confidence bands cover the finite-resolution target \(F_{h,0}\). A symmetric
second-order kernel gives a uniform \(O(h^2)\) difference from \(F_0\) under the
smoothness conditions in the Supplement.

Any threshold-indexed augmentation \(a_y(W)\) gives the unbiased signal
\begin{equation}
U_y(a)
=
a_y(W)+\frac{R}{\pi_M(W)}\{Z_y-a_y(W)\}.
\label{eq:pseudo}
\end{equation}
The signal satisfies \(E\{U_y(a)\mid W\}=m_{0,y}(W)\) and
\(E\{U_y(a)\mid X\}=q_{0,y}(X)\) for every \(a\). Setting \(a\equiv0\) gives
inverse-probability weighting. We call an \(X\)-based threshold regression the
reference-outcome anchor because it can be fitted with the verified outcomes
without using \(S\) in the augmentation. A regression using \((X,S)\) gives full
surrogate borrowing.

The augmentation changes variance without changing the target. Write
\(C_{yz}(W)=\operatorname{Cov}(Z_y,Z_z\mid W)\) and
\(\delta_y^a(W)=a_y(W)-m_{0,y}(W)\). Then
\[
\operatorname{Cov}\{U_y(a),U_z(a)\mid W\}
=
\frac{C_{yz}(W)}{\pi_M(W)}
+
\left\{\frac1{\pi_M(W)}-1\right\}
\delta_y^a(W)\delta_z^a(W).
\]
The first term is verified-outcome noise. The second is a positive-semidefinite
covariance contribution from augmentation error. Thus a surrogate can predict
\(Y\) well overall and still fail to reduce variance near \(x_0\). This identity
motivates an adaptive amount of borrowing.

\subsection{Adaptive surrogate borrowing}

Condition on the regression-fitting sample and suppress hats and fold indices.
Let \(q_y(X)\) be the anchor candidate and let \(m_y(W)\) use the surrogate. Their
difference \(D_y=m_y-q_y\) is the proposed correction; the unknown ideal
correction is \(A_y=m_{0,y}-q_y\). We use
\(a_{\lambda,y}=q_y+\lambda D_y\) with the same \(\lambda\) at every threshold.
The measure \(\nu\) specifies which thresholds matter for selection. We use
equal mass on a fixed grid in the primary analyses, although a tail-focused
analysis could place more mass in that region. For finite positive \(\nu\), the
integrated augmentation variance that depends on \(\lambda\) is
\[
\begin{aligned}
J_{M,h}(\lambda)
&=
E\{v_{M,h}\|A-\lambda D\|_\nu^2\}\\
&=
J_{M,h}(0)-2\lambda G_{M,h}+\lambda^2H_{M,h},\\
G_{M,h}
&=E\{v_{M,h}\langle A,D\rangle_\nu\},
\qquad
H_{M,h}=E\{v_{M,h}\|D\|_\nu^2\},
\end{aligned}
\]
where
\[
v_{M,h}(W)
=
\frac{h^dK_h^2(X-x_0)}{\mu_h^2}
\left\{\frac1{\pi_M(W)}-1\right\}.
\]
Here \(H_{M,h}\) measures the local size of the proposed correction and
\(G_{M,h}\) measures its alignment with the ideal correction. Full borrowing
reduces integrated augmentation variance exactly when
\(2G_{M,h}-H_{M,h}>0\).

When \(H_{M,h}>0\), the population choice is
\begin{equation}
\lambda_h^{\mathrm{or}}
=
\Pi_{[0,1]}
\left(
\frac{G_{M,h}}{H_{M,h}}
\right),
\label{eq:oracle-lambda}
\end{equation}
with \(\lambda_h^{\mathrm{or}}=0\) when \(H_{M,h}=0\). The dependence on \(M\)
is implicit. Missing-at-random verification gives an observable representation
of \(G_{M,h}\), even though \(A\) is unknown:
\[
G_{M,h}
=
E\left[
v_{M,h}
\left\langle
D,\frac{R}{\pi_M}(Z-q)
\right\rangle_\nu
\right].
\]
One common weight need not minimize variance at each threshold. We therefore
report lower- and upper-tail errors as well as integrated CDF error.

\subsection{Cross-fitting and simultaneous inference}

We keep weight selection separate from CDF evaluation. The sample is divided
into three folds: one fits the two regressions, another selects the weight, and
the third estimates the CDF. We rotate these assignments so that each
observation appears once in an evaluation fold. In rotation \(k\), the fitted
functions \(\widehat q_k\) and \(\widehat m_k\) are fixed before the selection
fold is used.

Let \(\mathbb P_{G,k}\) denote the selection-fold empirical mean and \(\widehat\mu_{h,G,k}=\mathbb P_{G,k}K_h(X-x_0)\). The empirical transfer weight is
\begin{equation}
\widehat v_{h,k}(W)
=
\frac{h^dK_h^2(X-x_0)}{\widehat\mu_{h,G,k}^{\,2}}
\left\{\frac1{\pi_M(W)}-1\right\}.
\label{eq:vhat}
\end{equation}
With \(\widehat D_k=\widehat m_k-\widehat q_k\), the selection fold computes
\[
\begin{aligned}
\widehat G_{h,k}
&=
\mathbb P_{G,k}
\left[
\widehat v_{h,k}
\left\langle
\widehat D_k,
\frac{R}{\pi_M}(Z-\widehat q_k)
\right\rangle_\nu
\right],\\
\widehat H_{h,k}
&=
\mathbb P_{G,k}
\{\widehat v_{h,k}\|\widehat D_k\|_\nu^2\},\\
\widehat\lambda_{h,k}
&=
\Pi_{[0,1]}
\left\{
\frac{\widehat G_{h,k}}
{\max(\widehat H_{h,k},\varepsilon_{H,M})}
\right\},
\end{aligned}
\]
where \(\varepsilon_{H,M}\downarrow0\).

On the evaluation fold, use the selected path
\[
\widehat a_{k,y}
=
\widehat q_{k,y}
+
\widehat\lambda_{h,k}
(\widehat m_{k,y}-\widehat q_{k,y})
\]
inside the augmented signal. Let \(\widehat F_{h,k}\) be the corresponding kernel ratio and let \(\rho_{M,k}\) be the evaluation-fold sample fraction. The combined raw estimator is
\begin{equation}
\widehat F_h^{\mathrm{raw}}(y\mid x_0)
=
\sum_{k=1}^3\rho_{M,k}\widehat F_{h,k}(y\mid x_0).
\label{eq:final-estimator}
\end{equation}
Conditional on the other two folds, the fitted functions and selected weight
are fixed when the evaluation-fold mean is formed. This separation is the basis
for the process limit and multiplier approximation developed below.

\begin{procedurebox}
\textbf{Algorithm 1: Local CDF estimation with adaptive surrogate borrowing}
\begin{enumerate}[leftmargin=1.6em,itemsep=0.2em,topsep=0.4em]
\item Fix \(x_0\), \(h\), the threshold grid, \(\nu\), and three folds.
\item Fit the \(X\)-based and surrogate-assisted threshold regressions on the fitting fold.
\item Estimate \(\widehat G_{h,k}\), \(\widehat H_{h,k}\), and \(\widehat\lambda_{h,k}\) on the selection fold.
\item Estimate the augmented local CDF on the evaluation fold.
\item Rotate the fold assignments and combine the evaluation-fold estimators using \eqref{eq:final-estimator}.
\item Construct multiplier critical values from the evaluation-fold influence
functions, and rearrange the point estimate for reporting and quantile
inversion.
\end{enumerate}
\end{procedurebox}

The raw estimate can leave \([0,1]\) or violate monotonicity. We report
\[
\widehat F_h^{\mathrm{mon}}
=
\mathcal R\{0\vee\widehat F_h^{\mathrm{raw}}\wedge1\},
\qquad
\widehat Q_h(\tau\mid x_0)
=
\inf\{y:\widehat F_h^{\mathrm{mon}}(y\mid x_0)\ge\tau\},
\]
where \(\mathcal R\) is increasing rearrangement. Inference for \(F_{h,0}\)
treats \(h\) as the inferential resolution and uses the raw-process band in
Section~\ref{sec:theory}. Inference for \(F_0\) also requires undersmoothing or
bias correction.
